% Part:sets-functions-relations
% Chapter: size-of-sets
% Section: enumerations

\documentclass[../../../include/open-logic-section]{subfiles}

\begin{document}

\olfileid{sfr}{siz}{enm}

\olsection{তালিকায়ন ও \usetoken{S}{enumerable} সেট}

\begin{editorial}
  এই অংশে সেটের তালিকায়ন আলোচনা করা হয়েছে; তালিকায়নকে $\PosInt$ থেকে সমাপতিত অপেক্ষক হিসেবে সংজ্ঞায়িত করা হয়েছে। গণিতের পূর্বপরিচয় অল্প এমন পাঠকদের জন্য আলোচনা ধীরে এগিয়েছে। \olref[enm-alt]{sec}-এ একটি বিকল্প, সংক্ষিপ্ততর সংস্করণ আছে; সেখানে তালিকায়নকে ভিন্নভাবে সংজ্ঞায়িত করা হয়েছে: $\Nat$-এর (বা তার কোনো প্রারম্ভিক খণ্ডের) সঙ্গে একৈক সমাপতিত অপেক্ষক হিসেবে।
\end{editorial}

\begin{explain}
আমরা ইতিমধ্যেই !!{element}s তালিকায় লিখে সেটের উদাহরণ দিয়েছি। একটি সেট অসীম হলেও কীভাবে এবং কখন তার !!{element}s তালিকায় লেখা যায়, এবার তা আরও সাধারণভাবে আলোচনা করা যাক।
\end{explain}

\begin{defn}[তালিকায়ন, অনানুষ্ঠানিকভাবে]
অনানুষ্ঠানিকভাবে, একটি সেট~$A$-এর \emph{তালিকায়ন} হলো~$A$-এর !!{element}s-এর একটি তালিকা (সম্ভবত অসীম), যাতে $A$-এর প্রত্যেক !!{element} কোনো সসীম ক্রমস্থানে আসে। $A$-এর কোনো তালিকায়ন থাকলে $A$-কে \emph{!!{enumerable}} বলা হয়।
\end{defn}

\begin{explain}
তালিকায়ন সম্পর্কে কয়েকটি কথা:
\begin{enumerate}
\item কেবল সেইসব তালিকাকেই তালিকায়ন ধরি, যাদের একটি শুরু আছে এবং যাদের প্রথম !!{element} ছাড়া প্রত্যেক !!{element}-এর অব্যবহিত আগে একটিমাত্র !!{element} আছে। অন্যভাবে বললে, তালিকার প্রথম !!{element} ও অন্য যেকোনো উপাদানের মধ্যে কেবল সসীমসংখ্যক উপাদান থাকে। বিশেষত এর অর্থ হলো, তালিকায়নের প্রত্যেক !!{element}-এর একটি সসীম ক্রমস্থান আছে: প্রথম !!{element}-এর ক্রমস্থান~$1$, দ্বিতীয়টির ক্রমস্থান~$2$, ইত্যাদি।
\item একই সেট~$A$-এর ভিন্ন তালিকায়নে !!{element}s ভিন্ন ক্রমে আসতে পারে: $4$, $1$, $25$, $16$,~$9$ যেমন প্রথম পাঁচটি বর্গসংখ্যার সেটের তালিকায়ন, তেমনি $1$, $4$, $9$, $16$,~$25$-ও তার তালিকায়ন।
\item পুনরাবৃত্তিযুক্ত তালিকায়নও তালিকায়ন: $1$, $1$, $2$, $2$, $3$, $3$,~\dots{} যে সেটের তালিকায়ন, $1$, $2$, $3$,~\dots{}-ও সেই একই সেটের তালিকায়ন।
\item তবে একটি নির্দিষ্ট তালিকায়ন বর্ণনা করার সময় ক্রম ও পুনরাবৃত্তি \emph{গুরুত্বপূর্ণ}। ধনাত্মক পূর্ণসংখ্যাগুলির তালিকায়ন $1$, $2$, $3$, $1$, \dots{} দিয়ে শুরু করতে পারি, কিন্তু প্রচলিতভাবে $1$, $2$, $3$, $4$,~\dots{} লিখলে বিন্যাসটি সহজে বোঝা যায়।
\item তালিকায়নের একটি শুরু থাকতে হবে: \dots, $3$, $2$, $1$ ধনাত্মক পূর্ণসংখ্যাগুলির তালিকায়ন নয়, কারণ এর কোনো প্রথম !!{element} নেই। অনানুষ্ঠানিক সংজ্ঞা থেকে এটি কীভাবে আসে বুঝতে নিজেকে জিজ্ঞাসা করো, ``এই তালিকায় 76 সংখ্যাটি কোন ক্রমস্থানে আছে?''
\item নিচের তালিকাটি ধনাত্মক পূর্ণসংখ্যাগুলির তালিকায়ন নয়: $1$, $3$, $5$, \dots, $2$, $4$, $6$, \dots\@ সমস্যাটি হলো, জোড় সংখ্যাগুলি সসীম ক্রমস্থানে না এসে $\infty + 1$, $\infty + 2$, $\infty +
  3$ ক্রমস্থানে আসে।
\item শূন্য সেট তালিকায়নযোগ্য: শূন্য তালিকাই তার তালিকায়ন!{}
\end{enumerate}
\end{explain}

\begin{prop}
  $A$-এর কোনো তালিকায়ন থাকলে তার পুনরাবৃত্তিবিহীন তালিকায়নও আছে।
\end{prop}

\begin{proof}
  ধরি, $A$-এর একটি তালিকায়ন $x_1$, $x_2$, \dots{} আছে, যেখানে প্রত্যেক $x_i$ হলো~$A$-এর একটি !!{element}। পুনরাবৃত্ত !!{element}s বাদ দিয়ে আমরা তালিকায়ন থেকে পুনরাবৃত্তি সরাতে পারি। যেমন, তালিকায়নটিকে নতুন একটি তালিকায়নে বদলাতে পারি এইভাবে: $x_i$ যদি~$A$-এর !!a{element} হয়, যা $x_1$, \dots, $x_{i-1}$-এর মধ্যে নেই, তবে তাকে তালিকায় রাখব; আর $x_i$ যদি ইতিমধ্যেই $x_1$, \dots,~$x_{i-1}$-এর মধ্যে থাকে, তবে তাকে তালিকা থেকে বাদ দেব।
\end{proof}

শেষ যুক্তিটি দেখায় যে তালিকায়ন ও !!{enumerable} সেটগুলিকে ভালোভাবে বুঝতে এবং তাদের সম্পর্কে বিভিন্ন ফল প্রমাণ করতে আমাদের আরও নির্দিষ্ট সংজ্ঞা দরকার। পরের সংজ্ঞাটি সেই প্রয়োজন মেটায়।

\begin{defn}[তালিকায়ন, আনুষ্ঠানিকভাবে]
একটি সেট $A \neq \emptyset$-এর \emph{তালিকায়ন} হলো যেকোনো !!{surjective} অপেক্ষক $f \colon \PosInt \to A$।
\end{defn}

\begin{explain}
আনুষ্ঠানিক সংজ্ঞা এবং সম্ভবত অসীম তালিকা ব্যবহার করা অনানুষ্ঠানিক সংজ্ঞাটি যে সমতুল্য, তা যাচাই করা যাক। প্রথমত, $\PosInt$ থেকে কোনো সেট~$A$-তে যেকোনো !!{surjective} অপেক্ষক~$A$-এর তালিকায়ন করে। এই ধরনের অপেক্ষক উপরের অনানুষ্ঠানিক সংজ্ঞা অনুযায়ী একটি তালিকায়ন নির্ধারণ করে: তালিকা $f(1)$, $f(2)$, $f(3)$, \dots। যেহেতু $f$ !!{surjective}, তাই~$A$-এর প্রত্যেক !!{element} অবশ্যই কোনো~$n \in \PosInt$-এর জন্য~$f(n)$-এর মান। সুতরাং $A$-এর প্রত্যেক !!{element} তালিকার কোনো সসীম ক্রমস্থানে আসে। অপেক্ষকটি !!{injective} না-ও হতে পারে, তাই তালিকায় পুনরাবৃত্তি থাকতে পারে; কিন্তু আগেই যেমন বলেছি, তাতে অসুবিধা নেই।

অন্যদিকে,~$A$-এর সকল !!{element}s-এর একটি তালিকায়ন দেওয়া থাকলে, $f(n)$-কে তালিকার $n$-তম !!{element} হিসেবে নিয়ে, অথবা $n$-তম !!{element} না থাকলে তালিকার শেষ !!{element} হিসেবে নিয়ে, আমরা !!a{surjective} অপেক্ষক $f\colon \PosInt \to A$ সংজ্ঞায়িত করতে পারি। কেবল $A$ শূন্য হলে, এবং সেই কারণে তালিকা শূন্য হলে, এই পদ্ধতি !!a{surjective} অপেক্ষক দেয় না। সুতরাং প্রত্যেক অশূন্য তালিকা !!a{surjective} অপেক্ষক $f\colon \PosInt \to A$ নির্ধারণ করে।
\end{explain}

\begin{defn}
  \ollabel{defn:enumerable}
  একটি সেট~$A$ !!{enumerable} হয় যদি এবং কেবল যদি সেটি শূন্য হয় অথবা তার একটি তালিকায়ন থাকে।
\end{defn}

\begin{ex}
ধনাত্মক পূর্ণসংখ্যাগুলির ($\PosInt$) তালিকায়নকারী একটি অপেক্ষক হলো $f(n) = n$ দ্বারা নির্ধারিত অভেদ অপেক্ষক। স্বাভাবিক সংখ্যা $\Nat$-এর তালিকায়নকারী একটি অপেক্ষক হলো $g(n) = n - 1$।
\end{ex}

\begin{ex}
নিচের সংজ্ঞা দ্বারা নির্ধারিত অপেক্ষক $f\colon \PosInt \to \PosInt$ ও $g \colon \PosInt \to
\PosInt$ যথাক্রমে জোড় ধনাত্মক পূর্ণসংখ্যা ও বিজোড় ধনাত্মক পূর্ণসংখ্যাগুলির তালিকায়ন করে:
\begin{align*}
f(n) & = 2n \text{ এবং}\\
g(n) & = 2n - 1
\end{align*}
তবে কোনো অপেক্ষকই $\PosInt$-এর তালিকায়ন নয়, কারণ কোনোটিই !!{surjective} নয়।
\end{ex}

\begin{prob}
ধনাত্মক বর্গসংখ্যা $1$, $4$, $9$, $16$, \dots{}-এর একটি তালিকায়ন সংজ্ঞায়িত করো।
\end{prob}

\begin{ex}
অপেক্ষক $f(n) = (-1)^{n} \lceil \frac{(n-1)}{2}\rceil$ পূর্ণসংখ্যার সেট~$\Int$-এর তালিকায়ন করে (এখানে $\lceil x \rceil$ হলো \emph{ঊর্ধ্ব পূর্ণাংশ} অপেক্ষক, যা $x$-কে উপরের দিকে নিকটতম পূর্ণসংখ্যায় নিয়ে যায়)। লক্ষ করো, $f$ ধনাত্মক ও ঋণাত্মক পূর্ণসংখ্যার মধ্যে এদিক-ওদিক ``লাফিয়ে'' কীভাবে $\Int$-এর মানগুলি তৈরি করে:
\[
\begin{array}{c c c c c c c c}
f(1) & f(2) & f(3) & f(4) & f(5) & f(6) & f(7) & \dots \\ \\
- \lceil \tfrac{0}{2} \rceil & \lceil \tfrac{1}{2}\rceil & - \lceil \tfrac{2}{2} \rceil & \lceil \tfrac{3}{2} \rceil & - \lceil \tfrac{4}{2} \rceil  & \lceil \tfrac{5}{2}
\rceil & - \lceil \tfrac{6}{2} \rceil & \dots \\ \\
0 & 1 & -1 & 2 & -2 & 3 & \dots
\end{array}
\]
$f$-কে নিচের মতো ক্ষেত্রভেদে সংজ্ঞায়িত বলেও ভাবতে পারো:
\[
f(n) = \begin{cases}
  0 & \text{যদি $n = 1$ হয়}\\
  n/2 & \text{যদি $n$ জোড় হয়}\\
  -(n-1)/2 & \text{যদি $n$ বিজোড় এবং $>1$ হয়}
  \end{cases}
\]
\end{ex}

\begin{prob}
  দেখাও যে $A$ ও $B$ !!{enumerable} হলে $A \cup B$-ও তাই। এর জন্য ধরি, !!{surjective} অপেক্ষক $f\colon \PosInt \to
  A$ ও $g\colon \PosInt \to B$ আছে। !!a{surjective} অপেক্ষক~$h\colon \PosInt \to A \cup B$ সংজ্ঞায়িত করো এবং সেটি !!{surjective} প্রমাণ করো। $A$ বা~$B = \emptyset$ হওয়ার ক্ষেত্রগুলিও বিবেচনা করো।
\end{prob}

\begin{prob}
  দেখাও যে $B \subseteq A$ এবং $A$ !!{enumerable} হলে~$B$-ও তাই। এর জন্য ধরি, !!a{surjective} অপেক্ষক $f\colon \PosInt \to
  A$ আছে। !!a{surjective} অপেক্ষক~$g\colon \PosInt \to B$ সংজ্ঞায়িত করো এবং সেটি !!{surjective} প্রমাণ করো। $B = \emptyset$ হলে কী হয়?
\end{prob}

\begin{prob}
  $n$-এর উপর গাণিতিক আরোহের সাহায্যে দেখাও যে $A_1$, $A_2$, \dots, $A_n$ সকলেই !!{enumerable} হলে $A_1 \cup \dots \cup A_n$-ও তাই। দুটি সেট $A$ ও~$B$ !!{enumerable} হলে~$A \cup
  B$-ও তাই—এই ফলটি তুমি ধরে নিতে পারো।
\end{prob}

একটি সেটের !!{element}s তালিকায় লেখার সময় $1$-তম !!{element} থেকে গোনা শুরু করা হয়তো বেশি স্বাভাবিক, কিন্তু গণিতবিদেরা গণনার জন্য স্বাভাবিক সংখ্যা~$\Nat$ ব্যবহার করতে পছন্দ করেন। তাঁরা তালিকার $0$-তম, $1$-তম, $2$-তম ইত্যাদি !!{element}s-এর কথা বলেন। সেই অনুযায়ী আমরা তালিকায়নকে $\Nat$ থেকে~$A$-তে !!a{surjective} অপেক্ষক হিসেবেও সংজ্ঞায়িত করতে পারি। অবশ্য দুটি সংজ্ঞা সমতুল্য।

\begin{prop}\ollabel{prop:enum-shift}
  !!a{surjection} $f\colon \PosInt \to A$ থাকে যদি এবং কেবল যদি !!a{surjection} $g\colon \Nat \to A$ থাকে।
\end{prop}

\begin{proof}
  !!a{surjection} $f\colon \PosInt \to A$ দেওয়া থাকলে, সকল $n \in \Nat$-এর জন্য $g(n) =
  f(n+1)$ সংজ্ঞায়িত করতে পারি। সহজেই দেখা যায় যে $g\colon \Nat
  \to A$ !!{surjective}। বিপরীতভাবে, !!a{surjection} $g\colon
  \Nat \to A$ দেওয়া থাকলে $f(n) = g(n-1)$ সংজ্ঞায়িত করো।
\end{proof}

এর থেকে নিচের ফলটি পাওয়া যায়:

\begin{cor}\ollabel{cor:enum-nat}
একটি সেট $A$ !!{enumerable} হয় যদি এবং কেবল যদি সেটি শূন্য হয় অথবা !!a{surjective} অপেক্ষক $f\colon \Nat \to A$ থাকে।
\end{cor}

আমরা উপরে আলোচনা করেছি যে একটি সেট~$A$-এর !!{element}s-এর তালিকাকে পুনরাবৃত্তিবিহীন তালিকায় বদলানো যায়। তালিকায়নের ক্ষেত্রেও এটি সত্য, তবে নিখুঁতভাবে বক্তব্যটি প্রকাশ করা ও প্রমাণ করা কিছুটা কঠিন। যেকোনো অপেক্ষক $f\colon \PosInt \to A$-কে সকল $n \in
\PosInt$-এর জন্য সংজ্ঞায়িত হতে হবে। যদি~$A$-তে কেবল সসীমসংখ্যক !!{element}s থাকে, তবে স্পষ্টতই~$\PosInt$-এর অসীমসংখ্যক !!{element}s-এর উপর সংজ্ঞায়িত এমন অপেক্ষক থাকতে পারে না, যা~$A$-এর সকল !!{element}s-কে মান হিসেবে নেয় অথচ কখনো একই মান দুবার নেয় না। সেই ক্ষেত্রে, অর্থাৎ পুনরাবৃত্তিবিহীন তালিকাটি সসীম হলে, আমাদের~$f$-এর জন্য একটি ভিন্ন সংজ্ঞাক্ষেত্র বেছে নিতে হবে, যাতে কেবল সসীমসংখ্যক !!{element}s আছে। পুনরাবৃত্তি না-থাকার অর্থ $f$-কে !!{injective} হতে হবে। যেহেতু সেটি একই সঙ্গে !!{surjective}, তাই আমরা কোনো সসীম সেট $\{1, \dots, n\}$ বা $\PosInt$ এবং~$A$-এর মধ্যে !!a{bijection} খুঁজছি।

\begin{prop}\ollabel{prop:enum-bij}
$f\colon \PosInt \to A$ যদি !!{surjective} হয় (অর্থাৎ~$A$-এর একটি তালিকায়ন হয়), তবে !!a{bijection} $g\colon Z \to A$ আছে, যেখানে $Z$ হয়~$\PosInt$, নয়তো কোনো~$n \in \PosInt$-এর জন্য $\{1, \dots, n\}$।
\end{prop}

\begin{proof}
  পূর্ববর্তী মানের সাহায্যে ধাপে ধাপে অপেক্ষক $g$ সংজ্ঞায়িত করি: ধরি, $g(1) = f(1)$। $g(i)$ ইতিমধ্যেই সংজ্ঞায়িত হয়ে থাকলে, $f(1)$, $f(2)$, \dots{}-এর যে প্রথম মানটি $g(1)$, \dots, $g(i)$-এর মধ্যে এখনও নেই, তেমন মান থাকলে সেটিই $g(i+1)$ নিই। $A$-তে ঠিক $n$টি !!{element}s থাকলে $g(1)$, \dots, $g(n)$ সকলেই সংজ্ঞায়িত হয়; ফলে আমরা একটি অপেক্ষক $g\colon \{1, \dots, n\}
  \to A$ সংজ্ঞায়িত করেছি। $A$-তে অসীমসংখ্যক !!{element}s থাকলে, যেকোনো $i$-এর জন্য তালিকায়ন $f(1)$, $f(2)$, \dots{}-এ~$A$-এর এমন !!a{element} অবশ্যই থাকে, যা $g(1)$, \dots, $g(i)$-এর মধ্যে এখনও নেই। এই ক্ষেত্রে আমরা একটি অপেক্ষক $g\colon \PosInt \to A$ সংজ্ঞায়িত করেছি।

  অপেক্ষক $g$ !!{surjective}, কারণ~$A$-এর যেকোনো উপাদান $f(1)$, $f(2)$, \dots{}-এর মধ্যে আছে (যেহেতু $f$ !!{surjective}); তাই সেটি অবশেষে কোনো~$i$-এর জন্য~$g(i)$-এর মান হবে। এটি !!{injective}-ও, কারণ $g(j) = g(i)$ হয় এমন $j < i$ থাকলে, $g(i)$ ইতিমধ্যেই $g(1)$, \dots, $g(i-1)$-এর মধ্যে থাকত, যা~$g$-এর সংজ্ঞার বিরোধী।
\end{proof}

\begin{cor}\ollabel{cor:enum-nat-bij}
একটি সেট $A$ !!{enumerable} হয় যদি এবং কেবল যদি সেটি শূন্য হয় অথবা !!a{bijection} $f\colon N \to A$ থাকে, যেখানে হয় $N = \Nat$, নয়তো কোনো $n \in \Nat$-এর জন্য $N = \{0, \dots, n\}$।
\end{cor}

\begin{proof}
$A$ !!{enumerable} হয় যদি এবং কেবল যদি $A$ শূন্য হয় অথবা !!a{surjective} $f\colon \PosInt \to A$ থাকে। \olref{prop:enum-bij} অনুসারে, দ্বিতীয় শর্তটি পূরণ হয় যদি এবং কেবল যদি !!a{bijective} অপেক্ষক~$f\colon Z \to A$ থাকে, যেখানে $Z =
\PosInt$ অথবা কোনো $n \in \PosInt$-এর জন্য $Z = \{1, \dots, n\}$। \olref{prop:enum-shift}-এর প্রমাণের একই যুক্তি অনুসারে, আবার এটি হয় যদি এবং কেবল যদি !!a{bijection} $g\colon N \to A$ থাকে, যেখানে হয় $N = \Nat$, নয়তো $N = \{0, \dots, n-1\}$।
\end{proof}

\begin{prob}
  \olref[sfr][siz][enm]{defn:enumerable} অনুসারে, একটি সেট $A$ তালিকায়নযোগ্য হয় যদি এবং কেবল যদি $A = \emptyset$ হয় অথবা !!a{surjective} $f\colon
  \PosInt \to A$ থাকে। ``!!{enumerable} সেট'' কথাটি নিচের নির্দিষ্ট সংজ্ঞা দিয়েও প্রকাশ করা যায়: একটি সেট তালিকায়নযোগ্য হয় যদি এবং কেবল যদি !!a{injective} অপেক্ষক $g\colon A \to \PosInt$ থাকে। সংজ্ঞা দুটি সমতুল্য দেখাও, অর্থাৎ দেখাও যে !!a{injective} অপেক্ষক $g\colon A \to \PosInt$ থাকে যদি এবং কেবল যদি হয় $A = \emptyset$, নয়তো !!a{surjective} $f\colon \PosInt \to A$ থাকে।
\end{prob}

\end{document}
